Convergence speed is strongly influenced by how anisotropic the curvature of $f$ is.
A standard way to quantify this is the \emph{condition number} of the Hessian.

\begin{definition}[Condition number]
If $f$ is twice differentiable and $\Hess f(x)$ is positive definite, define the (local) condition number
\begin{equation}
  \kappa(x) = \frac{\lambda_{\max}(\Hess f(x))}{\lambda_{\min}(\Hess f(x))}.
  \label{eq:condition-number-local}
\end{equation}
If $mI \preceq \Hess f(x) \preceq MI$ holds over a region of interest, then $\kappa(x)\le M/m$ on that region and we refer to $\kappa = M/m$ as a global condition number bound.
\end{definition}

\begin{example}[Bounding $\kappa$ for ridge logistic regression]
For $s(\theta)$ in \cref{ex:logistic-reg}, \eqref{eq:logistic-hess-bound}--\eqref{eq:logistic-ridge-hess} imply
\[
  mI \preceq \Hess s(\theta) \preceq \Bigl(m + \tfrac14 \lambda_{\max}(X^\top X)\Bigr) I,
\]
so we can take
\[
  M = m + \tfrac14 \lambda_{\max}(X^\top X),
  \qquad
  \kappa \le \frac{M}{m}.
\]
\end{example}

Geometrically, think of the level sets of $f$ near the optimum as ellipses.
When $\kappa$ is large, the ellipses are very elongated, so a single step size cannot simultaneously be ``right'' in both the steep and flat directions.
Vanilla gradient descent can therefore zig-zag and make slow progress; preconditioning (or second-order methods) aims to reduce the effective condition number.

\begin{figure}[t]
  \centering
  \begin{tikzpicture}[x=0.70cm,y=0.70cm,>=Stealth]
  % Level sets for f(x,y)=x^2/9 + y^2.
  \foreach \c/\col in {4/35,2/30,1/25,0.5/20,0.25/15} {
    \pgfmathsetmacro{\a}{3*sqrt(\c)}
    \pgfmathsetmacro{\b}{sqrt(\c)}
    \draw[stat221Gray!\col,line width=0.35pt] (0,0) ellipse ({\a} and {\b});
  }

  % Axes.
  \draw[->,stat221Gray!70,line width=0.45pt] (-6.4,0) -- (6.6,0);
  \draw[->,stat221Gray!70,line width=0.45pt] (0,-2.4) -- (0,2.6);
  \node[stat221Gray!70,font=\small,anchor=west] at (6.6,-0.10) {shallow};
  \node[stat221Gray!70,font=\small,anchor=south] at (-0.10,2.6) {steep};

  % Small step size (many small moves).
  \draw[stat221Red,thick,densely dotted,line join=round]
    (6.0000,2.0000) -- (5.8667,1.6000) -- (5.7363,1.2800) -- (5.6088,1.0240) -- (5.4842,0.8192)
    -- (5.3623,0.6554) -- (5.2431,0.5243) -- (5.1266,0.4194) -- (5.0127,0.3355) -- (4.9013,0.2684)
    -- (4.7924,0.2147) -- (4.6859,0.1718) -- (4.5818,0.1374) -- (4.4800,0.1100) -- (4.3804,0.0880)
    -- (4.2831,0.0704);
  \foreach \p in {(6.0000,2.0000),(5.8667,1.6000),(5.7363,1.2800),(5.6088,1.0240),(5.4842,0.8192),
    (5.3623,0.6554),(5.2431,0.5243),(5.1266,0.4194),(5.0127,0.3355),(4.9013,0.2684),
    (4.7924,0.2147),(4.6859,0.1718),(4.5818,0.1374),(4.4800,0.1100),(4.3804,0.0880),(4.2831,0.0704)}{
    \fill[stat221Red] \p circle (1.35pt);
  }

  % Large step size (faster but oscillatory in the steep direction).
  \draw[stat221Blue,thick,line join=round]
    (6.0000,2.0000) -- (4.9333,-1.2000) -- (4.0563,0.7200) -- (3.3352,-0.4320) -- (2.7423,0.2592)
    -- (2.2547,-0.1555) -- (1.8539,0.0933) -- (1.5243,-0.0560) -- (1.2533,0.0336) -- (1.0305,-0.0202)
    -- (0.8473,0.0121);
  \foreach \p in {(6.0000,2.0000),(4.9333,-1.2000),(4.0563,0.7200),(3.3352,-0.4320),(2.7423,0.2592),
    (2.2547,-0.1555),(1.8539,0.0933),(1.5243,-0.0560),(1.2533,0.0336),(1.0305,-0.0202),(0.8473,0.0121)}{
    \fill[stat221Blue] \p circle (1.35pt);
  }

  % Compact legend.
  \matrix[
    anchor=north west,
    draw=stat221Gray!25,
    fill=white,
    fill opacity=0.95,
    text opacity=1,
    rounded corners,
    inner sep=2.5pt,
    column sep=0.22cm,
    row sep=0.10cm,
    font=\small,
  ] at (-6.35,2.45) {
    \node{\tikz[baseline=-0.6ex]{\draw[stat221Red,thick,densely dotted] (0,0) -- (0.8,0); \fill[stat221Red] (0.4,0) circle (1.1pt);}}; &
      \node[stat221Red,anchor=west]{small step size}; \\
    \node{\tikz[baseline=-0.6ex]{\draw[stat221Blue,thick] (0,0) -- (0.8,0); \fill[stat221Blue] (0.4,0) circle (1.1pt);}}; &
      \node[stat221Blue,anchor=west]{large step size}; \\
  };
\end{tikzpicture}

  \caption{Ill-conditioned geometry: elongated level sets lead to slow, zig-zagging gradient descent.
  Small step sizes are stable but can be very slow; larger step sizes can oscillate in the steep direction (and must be controlled, e.g.\ by backtracking).}
  \label{fig:ill-conditioned-gd}
\end{figure}
